% ==================================================================================================================================
% An appendix for hires_correction.tex, on the radiation field a cloud requires of the models.
% To follow the appendix on the arrangement of the material along the line of sight.
% ==================================================================================================================================

\section{The radiation field a cloud requires of the models}
\label{app:field}

A correction that reads the fitted temperature can be applied only where the models cover the observations.  The grid of
Sect.~\ref{sec:law} heats its cores with a single radiation field, and a cloud illuminated more strongly is warmer at every
surface density than any model of that column, so its pixels fall outside the calibration however finely the latter is binned.
This appendix gives the procedure that decides, for a cloud in hand, which fields the models must have.  It is meant to be applied
by the reader to a cloud of their own and not only read.

\subsection{The scaling}
\label{app:field:scaling}

Dust that is optically thin to the radiation heating it settles where absorption balances emission.  The absorbed power is
proportional to the strength of the field, which we write as $G_0$ in units of the local interstellar field, and the power a grain
of opacity $\propto \lambda^{-\beta}$ emits is proportional to $T^{4+\beta}$, so

\begin{equation}
   T \propto G_0^{1/(4+\beta)} ,
\label{eq:g0scaling}
\end{equation}

which for the $\beta = 2$ adopted throughout is $T \propto G_0^{1/6}$.  This is the standard result for grains that are optically
thin to the radiation heating them \citep{IvezicElitzur1997, Tielens2005}, and we use it only as a scaling, since the calibration
itself comes from the radiative transfer and not from it.  The exponent is small, and that is the whole difficulty: a factor of
two in temperature demands a factor of $2^6 = 64$ in the field, so a cloud only moderately warmer than the models lies far
outside them in $G_0$, while a grid computed at one field covers a narrow band of temperature.  It is also what makes the remedy
affordable, since a ladder that is long in $G_0$ is short in temperature.

\subsection{The procedure}
\label{app:field:procedure}

The scaling is not applied to a cloud as a whole.  A pixel is as warm as it is because of the field that reaches it and because of
the column that shields it, and only the first is to be inferred.  The two are separated by comparing like with like, at a fixed
surface density:

\begin{enumerate}
\item Take the surface density image of the cloud and the temperature image that the same fit produced, both of which the
      reconstruction writes.  Bin the pixels by surface density.
\item In each bin, take the range of the observed temperature, for instance its 1st and 99th percentiles, which excludes the few
      pixels that any map has at its extremes.
\item In the same bin, take the range of temperature that the models give at the reference field $G_{\rm ref}$ of the existing
      grid.  This is a property of the grid and not of the cloud, and it is read from the same table the correction uses.
\item The field required to reach an observed temperature in that bin is
      \begin{equation}
         G_0 = G_{\rm ref} \left( \frac{T_{\rm observed}}{T_{\rm model}} \right)^{4+\beta} ,
      \label{eq:g0needed}
      \end{equation}
      evaluated at the low end of the observed range against the low end of the models, and likewise at the high end.
\item The ladder must span from the smallest such value over all bins to the largest.
\end{enumerate}

Only the ratio of the two temperatures enters, so nothing depends on the absolute calibration of either, and the shielding
divides out because both are taken at the same surface density.  Inferring a radiation field from a dust temperature is of course
common, and Eq.~(\ref{eq:g0needed}) adds nothing to the physics of it; what the comparison at a fixed surface density supplies is
a way to ask the narrower question this method needs answered, which is not whether a cloud is warm but whether the models reach
it.  For the diffuse gas, where nothing shields the dust, the model
temperature is the unattenuated $T_0$ of Eq.~(\ref{eq:law}) and the procedure reduces to the quick form

\begin{equation}
   G_0 = G_{\rm ref} \left( \frac{T}{T_0} \right)^{6} = 10 \left( \frac{T}{22.45\,{\rm K}} \right)^{6} ,
\label{eq:g0quick}
\end{equation}

which is accurate enough to place a cloud on the ladder before any of its data are binned.

\subsection{The procedure applied to Aquila}
\label{app:field:aquila}

\begin{table}
\caption{The procedure of Sect.~\ref{app:field:procedure} applied to the field of Sect.~\ref{sec:aquila}.  The third column is the
range of temperature the models of the grid span in that bin of surface density, the fourth the range the observations span, and
the last the fraction of the observed pixels that falls inside the models.}
\label{tab:g0aquila}
\centering
\begin{tabular}{lcccc}
\hline\hline
\noalign{\smallskip}
$\log_{10}\Sig$ & pixels & models (K) & Aquila (K) & inside \\
\noalign{\smallskip}
\hline
\noalign{\smallskip}
21.4 -- 21.6 &  1\,600 & 14.9 -- 17.3 & 15.9 -- 17.2 & 99.5{\%} \\
21.6 -- 21.8 & 26\,000 & 13.7 -- 15.9 & 14.8 -- 23.6 & 57.9{\%} \\
21.8 -- 22.0 & 78\,000 & 12.5 -- 14.7 & 14.1 -- 24.0 & 15.5{\%} \\
22.0 -- 22.2 & 88\,000 & 11.7 -- 13.7 & 13.2 -- 17.9 & 12.4{\%} \\
22.2 -- 22.4 & 66\,000 & 10.9 -- 11.9 & 12.6 -- 16.4 &  0{\%} \\
22.4 -- 22.6 & 42\,000 & 10.7 -- 11.1 & 12.0 -- 16.1 &  0{\%} \\
22.6 -- 22.8 & 27\,000 &  9.9 -- 10.5 & 11.6 -- 16.4 &  0{\%} \\
22.8 -- 23.0 & 15\,000 &  9.1 --  9.1 & 11.0 -- 15.9 &  0{\%} \\
above 23.0   &  9\,000 & no models    & 10.8 -- 16.3 & --      \\
\noalign{\smallskip}
\hline
\end{tabular}
\end{table}

Table~\ref{tab:g0aquila} carries the four steps out on the field of Sect.~\ref{sec:aquila}.  Over the whole field 17{\%} of the
pixels lie within the band of temperature that the models span at their surface density, and above $2{\times}10^{22}$\,cm$^{-2}$
none do: the dust of Aquila is one to five kelvin warmer than any model of the same column.  Inserting the ratios of the fourth
column to the third into Eq.~(\ref{eq:g0needed}) gives fields between about ten and about sixty, the upper end being reached in
the neighborhood of W40.  The grid of this paper, which has $G_0 = 10$, therefore covers the coolest of that field and nothing
else.

That table is the diagnostic to run before applying a temperature-dependent correction to any field.  It costs nothing, it uses
only images the reconstruction has already written, and it says at once whether the calibration covers the data or would be
extrapolating beyond them.

\subsection{Other clouds}
\label{app:field:others}

\begin{table}
\caption{The field required to reach the dust temperatures usually reported for the cold dust of each region, from the quick form
of Eq.~(\ref{eq:g0quick}).  These temperatures are indicative and serve only to place each cloud on the ladder; the procedure of
Sect.~\ref{app:field:procedure}, applied to the images of a particular field, supersedes them.}
\label{tab:g0}
\centering
\begin{tabular}{lcc}
\hline\hline
\noalign{\smallskip}
cloud & dust $T$ (K) & $G_0$ \\
\noalign{\smallskip}
\hline
\noalign{\smallskip}
Taurus              & 11 -- 16 & 0.14 -- 1.3 \\
Lupus, Scorpius     & 12 -- 20 & 0.23 -- 5.0 \\
Ophiuchus           & 12 -- 28 & 0.23 -- 38 \\
Aquila, with W40    & 11 -- 30 & 0.14 -- 57 \\
Cygnus~X            & 13 -- 32 & 0.38 -- 84 \\
Orion~A, with OMC-1 & 14 -- 45 & 0.59 -- 649 \\
\noalign{\smallskip}
\hline
\end{tabular}
\end{table}

Table~\ref{tab:g0} places six regions on the ladder with the quick form.  \draftnote{The temperature ranges of this table are
indicative values recalled from the published \emph{Herschel} temperature maps of these regions.  They must be replaced by values
read from the specific maps that will be used, with their references, before the paper is submitted.}  The lowest entries should
be read with the caution of Sect.~\ref{app:field:procedure}: a cloud whose densest material is at 11\,K does not require a field
of 0.14, because that material is cold through shielding and the same grid reaches it at a higher surface density.  What the low
end of the ladder has to reach is the coldest \emph{diffuse} gas, which for these regions implies $G_0$ between 0.7 and 1.3.

\subsection{The ladder}
\label{app:field:ladder}

Two rungs are fixed from outside: the interstellar field of the models cannot be scaled below $G_0 = 0.977$, and $G_0 = 10$ is the
grid already computed.  At a fixed surface density the models span about 17{\%} in temperature, because they differ in the
concentration and the size of their cores, so grids whose surface temperatures differ by less than that overlap and leave no gap
in the calibration.  Spacing the remaining rungs by the constant factor that joins the two fixed ones, $2.171$ in $G_0$ and
13.8{\%} in temperature, gives

\begin{equation}
   G_0 = 0.98,\ 2.12,\ 4.61,\ 10,\ 21.7,\ 47.1,\ 102,\ 222,\ 483 ,
\label{eq:g0ladder}
\end{equation}

which covers unattenuated temperatures from 15.2 to 42.8\,K in steps smaller than the spread the models themselves show.  It
reaches every region of Table~\ref{tab:g0}, the hottest gas of OMC-1 included.  Because the exponent of Eq.~(\ref{eq:g0scaling})
is $1/6$, those nine grids span a factor of 494 in the radiation field and only a factor of 2.8 in temperature.

The nodes need not all be recomputed at every field.  What a calibration requires is coverage in surface density, in
concentration and in the contrast between resolutions, and a subset of the lattice spanning the full range of each preserves it:
339 nodes at each of the nine fields is close to three times the cost of the single grid already computed, and that is the whole
price of making the correction applicable to clouds other than the coldest.
